\SourceLeaf{131}%
ubi $n$ est integer realis positivus, $A,B,C$ etc. integri reales vel imaginarii, $m$ autem integer complexus: vocabimus hic quoque \emph{radicem} congruentiae $X\equiv0\pmod m$ quemlibet integrum, qui pro $x$ substitutus ipsi $X$ valorem per modulum $m$ divisibilem conciliat. Solutiones per radices secundum modulum congruas non spectabimus tamquam diversas.

Quoties modulus est numerus primus, talis congruentia ordinis $n$ hic quoque plures quam $n$ solutiones diversas admittere non potest. Denotante $\alpha$ integrum quemvis determinatum (complexum), $X$ adiumento divisionis per $x-\alpha$ indefinite ad formam $X=(x-\alpha)X'+h$ reduci potest, ita ut $h$ fiat integer determinatus atque $X'$ functio ordinis $n-1$ cum coëfficientibus integris. Iam quoties $\alpha$ est radix congruentiae $X\equiv0\pmod m$, manifesto $h$ divisibilis erit per $m$, sive habebitur indefinite $X\equiv(x-\alpha)X'\pmod m$.

Perinde si denotante $\beta$ integrum determinatum, $X'$ ad formam $(x-\beta)X''+h'$ reducitur, $X''$ erit functio ordinis $n-2$ cum coëfficientibus integris. Si vero $\beta$ supponitur esse radix congruentiae $X\equiv0$, etiam satisfacere debet huic $(\beta-\alpha)X'\equiv0$, nec non huic $X'\equiv0$, siquidem radices $\alpha,\beta$ sunt incongruae, unde colligimus, etiam $h'$ per $m$ divisibilem esse debere, sive indefinite $X\equiv(x-\alpha)(x-\beta)X''\pmod m$.

Simili modo accedente radice tertia $\gamma$ prioribus incongrua, habebimus indefinite $X\equiv(x-\alpha)(x-\beta)(x-\gamma)X'''$, ita ut $X'''$ sit functio ordinis $n-3$ cum coëfficientibus integris. Eodem modo ulterius procedere licet, patetque simul, coëfficientem termini altissimi in singulis functionibus esse $=A$, quem per $m$ non divisibilem esse supponere licet, alioquin enim congruentia $X\equiv0$ essentialiter ad ordinem inferiorem referenda esset. Quoties itaque adsunt $n$ radices incongruae, puta $\alpha,\beta,\gamma\dotsfour\nu$, habebimus indefinite
\[
X\equiv A(x-\alpha)(x-\beta)(x-\gamma)\dotsfour(x-\nu)\pmod m
\]
quapropter substitutio novi valoris singulis $\alpha,\beta,\gamma\dotsfour\nu$ incongrui certo ipsi $X$ valorem per $m$ non divisibilem conciliaret, unde theorematis veritas sponte sequitur.

Ceterum haec demonstratio essentialiter convenit cum ea, quam in \emph{Disq. Ar.} art.~43 tradidimus, et cuius singula momenta pro numeris complexis perinde valent ac pro realibus.

\SourceLeaf{132}%
\Article{51}
Quae in Sectione tertia \emph{Disquisitionum Arithmeticarum} circa residua potestatum tradita sunt, ad maximam partem, levibus mutationibus adhibitis, etiam in arithmetica numerorum complexorum valent: quinadeo demonstrationes theorematum plerumque retineri possent. Ne tamen quid desit, theoremata principalia demonstrationibus concisis firmata proferemus, ubi semper subintelligendum est, modulum esse numerum primum.

\textsc{Theorema.} \emph{Denotante $k$ integrum per modulum $m$, cuius norma $=p$, non divisibilem, erit $k^{p-1}\equiv1\pmod m$.}

\emph{Demonstr.} Constituant $a,b,c$ etc. systema completum residuorum incongruorum pro modulo $m$, ita tamen, ut residuum per $m$ divisibile omissum sit, adeoque multitudo illorum numerorum, quorum complexum denotamus per $C$, sit $=p-1$. Sit porro $C'$ complexus productorum $ka,kb,kc$ etc. Ex his productis per hyp. nullum erit divisibile per $m$, quare singula habebunt residua congrua in complexu $C$, puta fieri poterit $ak\equiv a',bk\equiv b',ck\equiv c'$ etc. $(\operatorname{mod.}m)$, ita ut numeri $a',b',c'$ etc. ipsi in complexu $C$ inveniantur: denotemus complexum numerorum $a',b',c'$ etc. per $C''$. Sint $P,P',P''$ producta e singulis numeris complexuum $C,C',C''$ resp., sive
\[
\begin{aligned}
P&=abc\dotsfour\\
P'&=k^{p-1}abc\dotsfour=k^{p-1}P\\
P''&=a'b'c'\dotsfour
\end{aligned}
\]
Quum numeri complexus $C''$ deinceps congrui sint numeris complexus $C'$, erit $P''\equiv P'$ sive $P''\equiv k^{p-1}P$. At quum facile perspiciatur, binos quosvis numeros complexus $C''$ inter se incongruos, adeoque omnes inter se diversos esse, necessario numeri complexus $C''$ cum numeris complexus $C$ prorsus conveniunt, ordine tantummodo mutato, unde fit $P''=P$. Erit itaque $(k^{p-1}-1)P$ numerus per $m$ divisibilis, unde, quum $m$ sit numerus primus singulos factores ipsius $P$ non metiens, necessario $k^{p-1}-1$ per $m$ divisibilis esse debebit. Q.~E.~D.

\Article{52}
\textsc{Theorema.} \emph{Denotante $k$, ut in art. praec., integrum per modulum $m$ non divisibilem, atque $t$ exponentem minimum (praeter $0$), pro quo $k^t\equiv1\pmod m$, erit $t$ divisor cuiusvis alius exponentis $u$, pro quo $k^u\equiv1\pmod m$.}

\SourceLeaf{133}%
\emph{Demonstr.} Si $t$ non esset divisor ipsius $u$, sit $gt$ multiplum ipsius $u$ proxime maius quam $u$, adeoque $gt-u$ integer positivus minor quam $t$. Ex $k^t\equiv1,k^u\equiv1$, sequitur $0\equiv k^{gt}-k^u\equiv k^u(k^{gt-u}-1)$, adeoque $k^{gt-u}\equiv1$, i.~e. datur potestas ipsius $k$ cum exponente minori quam $t$ unitati congrua, contra hyp.

Tamquam corollarium hinc sequitur, $t$ certo metiri numerum $p-1$.

Numeros tales $k$, pro quibus $t=p-1$, etiam hic \emph{radices primitivas} pro. modulo $m$ vocabimus: quales revera adesse iam ostendemus.

\Article{53}
Resolvatur numerus $p-1$ in factores suos primos, ita ut habeatur
\[
p-1=a^\alpha b^\beta c^\gamma\dotsfour
\]
designantibus $a,b,c$ etc. numeros primos reales positivos inaequales. Sint $A,B,C$ etc. integri (complexi) per $m$ non divisibiles, atque resp. congruentiis
\[
x^{\frac{p-1}a}\equiv1,\quad x^{\frac{p-1}b}\equiv1,\quad x^{\frac{p-1}c}\equiv1\quad\text{etc.}
\]
secundum modulum $m$ \emph{non} satisfacientes, quales dari e theoremate art.~50 manifestum est. Denique sit $h$ congruus secundum modulum $m$ producto
\[
A^{\frac{p-1}{a^\alpha}}B^{\frac{p-1}{b^\beta}}C^{\frac{p-1}{c^\gamma}}\dotsfour
\]
Tunc dico, $h$ fore radicem primitivam.

\emph{Demonstr.} Denotando per $t$ exponentem infimae potestatis $h^t$ unitati congruae, erit, si $h$ non esset radix primitiva, $t$ submultiplum ipsius $p-1$, sive $\frac{p-1}t$ integer unitate maior. Manifesto hic integer factores suos primos reales inter hos $a,b,c$ etc. habebit: supponamus itaque, (quod licet), $\frac{p-1}t$ esse divisibilem per $a$, statuamusque $p-1=atu$. Erit itaque, propter $h^t\equiv1$, etiam $h^{tu}\equiv1$ sive
\[
A^{\frac{p-1}{a^\alpha}\pdot\frac{p-1}a}B^{\frac{p-1}{b^\beta}\pdot\frac{p-1}a}C^{\frac{p-1}{c^\gamma}\pdot\frac{p-1}a}\dotsfour\equiv1
\]
At manifesto $\frac{p-1}{ab^\beta}$ est integer, adeoque
\[
B^{\frac{p-1}{b^\beta}\pdot\frac{p-1}a}=(B^{p-1})^{\frac{p-1}{ab^\beta}}\equiv1
\]
perinde etiam
\SourceLeaf{134}%
$C^{\frac{p-1}{c^\gamma}\pdot\frac{p-1}a}\equiv1$, et sic porro; quapropter esse debet $A^{\frac{p-1}{a^\alpha}\pdot\frac{p-1}a}\equiv1$
Iam determinetur integer positivus $\lambda$ talis, ut fiat
\[
\lambda b^\beta c^\gamma\dotsfour\equiv1\pmod a
\]
quod fieri poterit, quum numerus primus $a$ ipsum $b^\beta c^\gamma\dotsfour$ non metiatur, statuaturque $\lambda b^\beta c^\gamma\dotsfour=1+a\mu$. Manifesto fit $A^{\lambda\pdot\frac{p-1}{a^\alpha}\pdot\frac{p-1}a}\equiv1$, sive, quoniam $\lambda\pdot\frac{p-1}{a^\alpha}\pdot\frac{p-1}a=(1+a\mu)\frac{p-1}a=(p-1)\mu+\frac{p-1}a$ habemus $A^{(p-1)\mu}\pdot A^{\frac{p-1}a}\equiv1$ atque hinc, quum sponte sit $A^{(p-1)\mu}\equiv1$, etiam $A^{\frac{p-1}a}\equiv1$, quod est contra hypothesin. Suppositio itaque, $t$ esse submultiplum ipsius $p-1$, consistere nequit, eritque adeo necessario $h$ radix primitiva.

\Article{54}
Denotante $h$ radicem primitivam pro modulo $m$, cuius norma $=p$, termini progressionis
\[
1,h,hh,h^3\dotsfour h^{p-2}
\]
inter se incongrui erunt, unde facile colligitur, quemlibet integrum non divisibilem per modulum uni ex istis congruum esse debere, sive illam seriem exhibere systema completum residuorum incongruorum exclusa cifra. Exponens eius potestatis, cui numerus datus congruus est, vocari potest huius \emph{index}, dum $h$ tamquam \emph{basis} consideratur. Ecce quaedam exempla, ubi cuivis indici residuum absolute minimum apposuimus.

\begin{center}\emph{Exemplum primum.}\end{center}
\[
m=5+4i,\quad p=41,\quad h=1+2i
\]
\begingroup\small
\[
\setlength{\arraycolsep}{2pt}
\begin{array}{r|r||r|r||r|r||r|r||r|r}
\text{Ind.}&\text{Residuum}&\text{Ind.}&\text{Residuum}&\text{Ind.}&\text{Residuum}&\text{Ind.}&\text{Residuum}&\text{Ind.}&\text{Residuum}\\\hline
0&+1&8&-4&16&-2+2i&24&+2i&32&+1+i\\
1&+1+2i&9&-3+i&17&-1+2i&25&-3i&33&+1+3i\\
2&+1-i&10&-i&18&+4i&26&+2+2i&34&+2\\
3&+3+i&11&+2-i&19&+1+3i&27&+2+i&35&-3\\
4&-2i&12&-1-i&20&-1&28&+4&36&+2-2i\\
5&+3i&13&+1-3i&21&-1-2i&29&+3-i&37&+1-2i\\
6&-2-2i&14&-2&22&-1+i&30&+i&38&-4i\\
7&-2-i&15&+3&23&-3-i&31&-2+i&39&-1-3i
\end{array}
\]
\endgroup

\SourceLeaf{135}%
\begin{center}\emph{Exemplum secundum.}\end{center}
\[
m=7,\quad p=49,\quad h=1+2i
\]
\begingroup\small
\[
\setlength{\arraycolsep}{2pt}
\begin{array}{r|r||r|r||r|r||r|r||r|r}
\text{Ind.}&\text{Residuum}&\text{Ind.}&\text{Residuum}&\text{Ind.}&\text{Residuum}&\text{Ind.}&\text{Residuum}&\text{Ind.}&\text{Residuum}\\\hline
0&+1&10&-1-i&20&+2i&30&+2-2i&40&+3\\
1&+1+2i&11&+1-3i&21&+3+2i&31&-1+2i&41&+3-i\\
2&-3-3i&12&-i&22&-1+i&32&+2&42&-2-2i\\
3&+3-2i&13&+2-i&23&-3-i&33&+2-3i&43&+2+i\\
4&-3i&14&-3+3i&24&-1&34&+1+i&44&-2i\\
5&-1-3i&15&-2-3i&25&-1-2i&35&-1+3i&45&-3-2i\\
6&-2+2i&16&-3&26&+3+3i&36&+i&46&+1-i\\
7&+1-2i&17&-3+i&27&-3+2i&37&-2+i&47&+3+i\\
8&-2&18&+2+2i&28&+3i&38&+3-3i&&\\
9&-2+3i&19&-2-i&29&+1+3i&39&+2+3i&&
\end{array}
\]
\endgroup

\Article{55}
Adiicimus circa radices primitivas et algorithmum indicum quasdam observationes, demonstrationibus propter facilitatem omissis.

I. Indices secundum modulum $p-1$ congrui in systemate dato residuis secundum modulum $m$ congruis respondent et vice versa.

II. Residua, quae respondent indicibus ad $p-1$ primis, etiam sunt radices primitivae et vice versa.

III. Si accepta radice primitiva $h$ pro basi, radicis alius primitivae $h'$ index est $t$, et vice versa $t'$ index ipsius $h$, dum $h'$ pro basi accipitur, erit $tt'\equiv1\pmod{p-1}$; et si iisdem positis indices cuiusdam alius numeri in his duobus systematibus resp. sunt $u,u'$, erit $tu'\equiv u,t'u\equiv u'\pmod{p-1}$.

IV. Dum numeri $1,1+i$ eorumque terni socii (tamquam nimis ieiuni) a modulis nobis considerandis excluduntur, restant numeri primi ii, quos in art.~34 tertio et quarto loco posuimus. Posteriorum normae erunt numeri primi reales formae $4n+1$; priorum normae autem quadrata numerorum primorum realium imparium: in utroque igitur casu $p-1$ per $4$ divisibilis est.

V. Denotando indicem numeri $-1$ per $u$, erit $2u\equiv0\pmod{p-1}$, adeoque vel $u\equiv0$, vel $u\equiv\frac12(p-1)$: at quum index $0$ respondeat residuo $+1$, index numeri $-1$ necessario debet esse $\frac12(p-1)$.

VI. Perinde denotando per $u$ indicem numeri $i$, erit $2u\equiv\frac12(p-1)\pmod{p-1}$, adeoque vel $u\equiv\frac14(p-1)$, vel $u\equiv\frac34(p-1)$. Sed hic ambiguitas ab electione radicis primitivae pendet. Scilicet si radice primitiva $h$ pro basi ac%
\SourceLeaf{136}%
cepta index numeri $i$ est $\frac14(p-1)$, index fiet $\frac34(p-1)$, dum pro basi accipitur $h^\mu$, designante $\mu$ integrum positivum formae $4n+3$ ad $p-1$ primum, e.~g. ipsum numerum $p-2$, et vice versa. Quare semissis altera radicum primitivarum conciliat numero $i$ indicem $\frac14(p-1)$, altera indicem $\frac34(p-1)$, manifestoque pro illis basibus $-i$ indicem $\frac34(p-1)$, pro his indicem $\frac14(p-1)$ habebit.

VII. Quoties modulus est numerus primus realis positivus formae $4n+3$, puta $=q$, adeoque $p=qq$, indices omnium numerorum realium per $q+1$ divisibiles erunt; denotante enim $t$ indicem numeri realis $k$, erit, propter $k^{q-1}\equiv\pmod q$, $(q-1)t\equiv0\pmod{qq-1}$, adeoque $\frac{t}{q+1}$ integer. Perinde indices numerorum pure imaginariorum ut $ki$ per $\frac12(q+1)$ divisibiles erunt. Patet itaque, radices primitivas pro talibus modulis inter solos numeros mixtos quaerendas esse.

VIII. Contra pro modulo $m$, qui est numerus primus complexus mixtus, (cuiusque proin norma $p$ est numerus primus realis formae $4n+1$), radices primitivae quaelibet etiam inter numeros reales eligi possunt, inter quos completum adeo systema residuorum incongruorum monstrare licet (art.~40). Manifesto autem quilibet numerus realis, qui est radix primitiva pro modulo complexo $m$, simul erit in arithmetica numerorum realium radix primitiva pro modulo $p$, et vice versa.

\Article{56}
Etiamsi theoria residuorum et non-residuorum quadraticorum in arithmetica numerorum complexorum sub ipsa theoria residuorum biquadraticorum contenta sit, tamen antequam ad hanc transeamus, illius theoremata palmaria hic seorsim proferemus: brevitatis vero caussa de solo casu principali, ubi modulus est numerus primus complexus (impar), hic loquemur.

Sit $m$ talis modulus, atque $p$ eius norma. Manifesto quivis integer (per $m$ non divisibilis, quod hic semper subintelligendum) quadrato secundum modulum $m$ congruus fieri vel potest vel non potest, prout illius index, radice aliqua primitiva pro basi accepta, par est vel impar; in casu priori ille integer residuum quadraticum ipsius $m$ dicetur, in posteriori non-residuum. Hinc concluditur, inter $p-1$ numeros qui systema completum residuorum incongruorum (per $m$ non divisibilium) exhibeant, semissem ad residua quadratica, semissem alteram ad non-residua quadratica referri. Cuivis vero alii numero extra illud systema idem
\SourceLeaf{137}%
character hoc respectu tribuendus est, quo gaudet numerus systematis illi congruus.

Porro ibinde sequitur, productum e duobus residuis quadraticis, nec non productum e duobus non-residuis esse residuum quadraticum, contra productum e residuo quadratico in non-residuum fieri non-residuum; et generaliter productum e quotcunque factoribus esse residuum quadraticum vel non-residuum, prout multitudo non-residuorum inter factores par sit vel impar.

Pro distinguendis residuis quadraticis a non-residuis statim se offert criterium generale sequens:

Numerus $k$ per modulum non divisibilis huius residuum vel non-residuum quadraticum est, prout habetur vel $k^{\frac12(p-1)}\equiv1$, vel $k^{\frac12(p-1)}\equiv-1\pmod m$.

Veritas huius theorematis statim inde sequitur, quod, accepta radice primitiva quacunque pro basi, index potestatis $k^{\frac12(p-1)}$ fit vel $\equiv0$ vel $\equiv\frac12(p-1)$, prout index numeri $k$ par est vel impar.

\Article{57}
Facile quidem est, pro modulo dato systema residuorum incongruorum completum in duas classes, puta residua et non-residua quadratica distinguere, quo pacto simul omnibus reliquis numeris classes suae sponte assignantur. At longe altioris indaginis est quaestio de criteriis ad distinguendum modulos eos, pro quibus numerus datus est residuum quadraticum, ab iis, pro quibus est non-residuum.

Quod quidem attinet ad unitates reales $+1$ et $-1$, hae in arithmetica numerorum complexorum sunt reapse quadrata, adeoque etiam residua quadratica pro \emph{quovis} modulo. Aeque facile e criterio art. praec. sequitur, numerum $i$ (et perinde $-i$) esse residuum quadraticum cuiusvis moduli, cuius norma $p$ sit formae $8n+1$, non-residuum vero cuiusvis moduli, cuius norma sit formae $8n+5$. Quum manifesto nihil intersit, utrum numerus $m$, an aliquis numerorum ipsi associatorum $im,-m,-im$ pro modulo adoptetur, supponere licebit, modulum esse associatorum primarium (art.~36, II), adeoque statuendo modulum $=a+bi$, esse $a$ imparem, $b$ parem. Quo pacto quum semper sit $aa\equiv1\pmod8$, $bb$ vero vel $\equiv0$ vel $\equiv4\pmod8$, prout $b$ sit pariter par vel impariter par, patet, numeros $+i$ et $-i$ in casu priori esse residua quadratica moduli, in posteriori non-residua.

\SourceLeaf{138}%
\Article{58}
Quum diiudicatio characteris numeri compositi, utrum sit residuum quadraticum an non-residuum, pendeat a characteribus factorum, manifesto sufficiet, si evolutionem criteriorum ad distinguendos modulos, pro quibus numerus datus $k$ sit residuum quadraticum, ab iis, pro quibus sit non-residuum, ad tales valores ipsius $k$ limitemus, qui sint numeri primi, insuperque inter associatos primarii. In qua investigatione \emph{inductio} protinus theoremata maxime elegantia suppeditat.

Incipiamus a numero $1+i$, qui invenitur esse residuum quadraticum modulorum
\[
\begin{gathered}
-1+2i,\ +3-2i,\ -5-2i,\ -1-6i,\ +5+4i,\ +5-4i,\\
-7,\ +7+2i,\ -5+6i,\quad\text{etc.}
\end{gathered}
\]
non-residuum quadraticum autem sequentium
\[
\begin{gathered}
-1-2i,\ -3,\ +3+2i,\ +1+4i,\ +1-4i,\ -5+2i,\\
-1+6i,\ +7-2i,\ -5-6i,\ -3+8i,\ -3-8i,\\
+5+8i,\ +5-8i,\ +9+4i,\ +9-4i\quad\text{etc.}
\end{gathered}
\]

Si hunc conspectum, in quo semper quaternis modulis associatis primarium apposuimus, attente examinamus, facile animadvertimus, modulos $a+bi$ in priori classe omnes esse tales, pro quibus $a+b$ fiat $\equiv+1\pmod8$, in posteriori vero tales, pro quibus $a+b\equiv-3\pmod8$. Manifesto hoc criterium, si loco moduli primarii $m$ adoptamus associatum $-m$, ita immutari debet, ut pro modulis prioris classis sit $a+b\equiv-1$, pro modulis posterioris $\equiv+3\pmod8$. Quare, siquidem inductio non fefellerit, generaliter, designante $a+bi$ numerum primum, in quo $a$ impar. $b$ par, $1+i$ fit eius residuum quadraticum vel non-residuum quadraticum, prout $a+b\equiv\pm1$, vel $\equiv\pm3\pmod8$.

Pro numero $-1-i$ eadem regula valet, quae pro $1+i$. Contra considerando $1-i$ tamquam productum ex $-i$ in $1+i$, manifestum est, numero $1-i$ eundem characterem competere, qui tribuendus sit ipsi $1+i$, quoties $b$ sit pariter par, oppositum autem, quoties $b$ sit impariter par, unde facile colligitur, $1-i$ esse residuum quadraticum numeri primi $a+bi$, quoties sit $a-b\equiv\pm1$, non-residuum autem, quoties habeatur $a-b\equiv\pm3\pmod8$, semper supponendo, $a$ esse imparem, $b$ parem.

Ceterum haec secunda propositio e priori etiam deduci potest adiumento theorematis generalioris, quod ita enunciamus:

\SourceLeaf{139}%
In theoria residuorum quadraticorum character numeri $\alpha+\beta i$ respectu moduli $a+bi$ idem est, qui numeri $\alpha-\beta i$ respectu moduli $a-bi$.

Demonstratio huius theorematis inde petitur, quod uterque modulus eandem normam $p$ habet, atque quoties $(\alpha+\beta i)^{\frac12(p-1)}-1$ per $a+bi$ divisibilis est, etiam $(\alpha-\beta i)^{\frac12(p-1)}-1$ per $a-bi$ divisibilis evadit, quoties autem $(\alpha+\beta i)^{\frac12(p-1)}+1$ per $a+bi$ divisibilis est, etiam $(\alpha-\beta i)^{\frac12(p-1)}+1$ per $a-bi$ divisibilis esse debet.

\Article{59}
Progrediamur ad numeros primos impares.

Numerum $-1+2i$ invenimus esse residuum quadraticum modulorum $+3+2i$, $+1-4i$, $-5+2i$, $-5-2i$, $-1-6i$, $+7-2i$, $-3+8i$, $+5+8i$, $+5-8i$, $+9+4i$ etc.

non-residuum autem modulorum $-1-2i$, $-3$, $+3-2i$, $+1+4i$, $-1+6i$, $+5+4i$, $+5-4i$, $-7$, $+7+2i$, $-5+6i$, $-5-6i$, $-3-8i$, $+9-4i$ etc.

Reducendo modulos prioris classis ad residua eorum absolute minima secundum modulum $-1+2i$, haec sola invenimus $+1$ et $-1$, puta $+3+2i\equiv-1$, $+1-4i\equiv-1$, $-5+2i\equiv+1$, $-5-2i\equiv-1$ etc.

Contra omnes moduli posterioris classis congrui inveniuntur secundum modulum $-1+2i$ vel ipsi $+i$, vel ipsi $-i$.

At numeri $+1,-1$ ipsi sunt residua quadratica moduli $-1+2i$, atque $+i$ et $-i$ eiusdem non-residua: quocirca, quatenus inductioni fidem habere licet, prodit theorema: Numerus $-1+2i$ est residuum vel non-residuum quadraticum numeri primi $a+bi$, prout hic est residuum vel non-residuum quadraticum ipsius $-1+2i$, siquidem $a+bi$ est primarius e quaternis associatis, vel potius, si $a$ est impar, $b$ par.

Ceterum ex hoc theoremate sponte sequuntur theoremata analoga circa numeros $+1-2i$, $-1-2i$, $+1+2i$.

\Article{60}
Instituendo similem inductionem circa numerum $-3$ vel $+3$, invenimus, utrumque esse residuum quadraticum modulorum $+3+2i$, $+3-2i$,
\SourceLeaf{140}%
$-1+6i$, $-1-6i$, $-7$, $-5+6i$, $-5-6i$, $-3+8i$, $-3-8i$, $+9+4i$, $+9-4i$ etc.

non-residuum vero horum $-1+2i$, $-1-2i$, $+1+4i$, $+1-4i$, $-5+2i$, $-5-2i$, $+5+4i$, $+5-4i$, $+7+2i$, $+7-2i$, $+5+8i$, $+5-8i$ etc.

Priores secundum modulum $3$ congrui sunt alicui ex his quatuor numeris $+1,-1,+i,-i$; posteriores autem alicui ex his $+1+i,+1-i,-1+i,-1-i$. Illi sunt ipsa residua quadratica numeri $3$, hi non-residua.

Docet itaque haec inductio, numerum primum $a+bi$, supponendo $a$ imparem, $b$ parem, ad numerum $-3$ (nec non ad $+3$) eandem relationem habere, quam hic habet ad illum, quatenus scilicet alter alterius residuum quadraticum sit aut non-residuum.

Extendendo similem inductionem ad alios numeros primos, ubique hanc elegantissimam reciprocitatis legem confirmatam invenimus, deferimurque ad theorema hocce fundamentale circa residua quadratica in arithmetica numerorum complexorum:

\emph{Denotantibus $a+bi,A+Bi$ numeros primos tales, ut $a,A$ sint impares, $b,B$ pares: erit vel uterque alterius residuum quadraticum, vel uterque alterius non-residuum.}

At non obstante summa theorematis simplicitate, ipsius demonstratio magnis difficultatibus premitur, quibus tamen hic non immoramur, quum theorema ipsum sit tantummodo casus specialis theorematis generalioris, summam theoriae residuorum biquadraticorum quasi exhaurientis. Ad hanc igitur iam transeamus.

\Article{61}
Quae in art.~2 prioris commentationis de notione residui et non-residui biquadratici prolata sunt, etiam ad arithmeticam numerorum complexorum extendimus, et perinde ut illic etiam hic disquisitionem ad modulos tales, qui sunt numeri primi, restringimus: simul plerumque tacite subintelligendum erit, modulum ita accipi, ut sit inter associatos primarius, puta $\equiv1$ secundum modulum $2+2i$, nec non numeros, de quorum charactere (quatenus sint residua biquadratica vel non-residua) agitur, per modulum non esse divisibiles.

Pro modulo itaque dato numeri per eum non divisibiles in tres classes dispertiri possent, quarum prima contineret residua biquadratica, secunda non-residua biquadratica ea, quae sunt residua quadratica, tertia non-residua quadratica.

